\section{Discussion and limitations}\label{sec:discussion}
\subsection{What the certificate contributes}
The results support a decision-oriented use of compact representations.
Authenticated old groups make the assigned-center contribution exact, and
switching bounds determine whether that information is sufficient to order
the two exact quantized costs. The 2/19 versus 19/19 ablation isolates the
value of switching information on one development workload. Natural
worsening updates with nonzero switching further show that the mechanism
can certify regressions rather than only confirm early Lloyd improvements.

This distinction matters when an approximate computation has high label
agreement but cannot establish the sign of a small aggregate change. Our
interval makes that uncertainty explicit. A valid unresolved request
receives further authenticated computation; high agreement alone never
authorizes an update. Conversely, an aggregate cost certificate does not
require every approximate label to be correct. The three reported outputs
therefore answer different questions and should remain separate when
selecting an operating profile.

The centroid-informed representation has modest agreement and label-coverage
advantages over PCA on the confirmation workloads, while both attain the
same final update coverage. Repeated MNIST timings are also close. The
principal contribution is the shared verified mechanism and its fallback
accounting. The evidence does not establish a generally superior predictor,
a necessary certificate rank of 79 or a universally best basis.

\subsection{Where the measured savings apply}
The completed Fashion control gives a measured Full/certificate cost ratio
of 2.388 for the implemented reusable exact Full path, including
preparation, proof work and replay. Repeated smaller controls support approximately twofold
savings on two fixed MNIST fits. This is useful evidence for release batches
with a fixed committed catalog and deployed reference. The total saving
depends on how many candidates share preprocessing and how often the
switching interval avoids complete fallback.

There are corresponding costs. Compact authentication produces more proofs
and requires more public verification than Full in the Fashion comparison.
Replacing the reference requires new group preprocessing. A known future
center bank can favor the exact affine alternative, while moments-only can
be sufficient for easily certified improvements. A deployment should select
among these contracts using expected candidate availability, reference
lifetime, disclosure permissions and measured total work. The current study
does not supply a universal speed ranking across these settings.

\subsection{Application breadth and empirical limits}
The sensor release batches extend the evaluation beyond related image
benchmarks and beyond candidates drawn only from a Lloyd trajectory. Their
11/12 decision coverage and high agreement identify a favorable additional
regime. Frozen neural text provides a more difficult case: every observed
regression remains unresolved before Full fallback, and assignment agreement
misses the working target. This failure is consistent with insufficiently
informative compact bounds for those updates; it does not invalidate a
conclusive certificate on another request.

The dataset coverage and number of independent fitted configurations remain
limited. HAR windows can be correlated within subjects, and documents can
share source or topic structure. Prefixes, reference ages, candidate rates
and basis variants reuse fitted models. Counts of requests therefore exceed
the number of independent learning trials. Ground-truth classification
performance and encoder correctness are outside the certificate, so the
results cannot establish that certified distortion improvement improves
application accuracy.

The largest reported native catalogs contain 16,384 records, but those
experiments do not measure proving at that scale. Primary timing comparisons
use one laptop, small center counts and a limited number of executions;
the Fashion headline has one execution. Conditional timing intervals
characterize fixed workloads, not population-level runtime variation.
The alternative-backend export also lacks an uninterrupted comparative
timing. These boundaries limit both deployment extrapolation and claims
against specialized matrix-product or GKR proving systems.

\subsection{Disclosure and deployment requirements}
The method targets authorized aggregate disclosure. The explicit probe
construction recovers group moments and may reveal singleton records.
Per-batch intervals and row corrections add observable information, even
though original rows and moments are circuit witnesses. Applications
requiring adaptive record privacy need a different output mechanism and
its own privacy analysis. Restricting queries is an authorization measure;
it is not a proof of non-reconstruction.

Accepted decisions additionally depend on authentic owner pins, sound
arguments, binding commitments, checked integer ranges, full catalog
coverage and an independent post-commit challenge. The local benchmark
checks these relations under simulated roles. Deployment requires
authenticated challenge delivery and setup governance; an independent
security audit remains outstanding. These requirements should be assessed
alongside performance rather than inferred from successful proof replay.

\subsection{Further evaluation}
The most direct next measurements are complete matched proof comparisons
on additional sensor and neural-feature fits, uninterrupted backend
controls, and release sequences that charge reference replacement. Rank
and refinement-budget studies should measure the joint effect on decision
coverage, preparation and query cost. Larger-center workloads would test
whether the compact switching relation retains its advantage. A total-only
output interface or proof aggregation could change disclosure and verifier
cost, but would require new implemented relations and corresponding
analysis. These directions are prospective work and are not counted as
completed evidence.

\section{Conclusion}\label{sec:conclusion}
We presented reusable zero-knowledge certificates for comparing candidate
prototype banks on a fixed committed feature catalog. Authenticated deployed
groups yield an exact assigned-center cost change; integer-safe switching
bounds and verified exact-row corrections complete a conservative interval
for the exact quantized objective. Unresolved requests use newly generated
exact Full proofs, with their cost retained in the service total.

The completed 1,024-record Fashion workload certifies all three worsening
updates in 42.534 minutes, compared with 101.575 minutes for reusable exact
Full. Frozen-profile numerical confirmation resolves all 60 image requests,
while HAR resolves 11/12 and neural text resolves 8/12. The latter leaves
all four observed regressions unresolved, identifying a concrete operating
limit. Together with the verification and reuse costs, these results
support a selective audit primitive for learning-system releases under
authorized aggregate disclosure. Its value is determined by decision
resolution and total work within that contract.
